% Revised from the 2026-11-19 single-execution raw measurements; totals include setup.
\section{Experimental Results}
\label{sec:experimental-results}
This section evaluates the performance of our AF-CPSI and the resulting
PJC protocol that combines it with homomorphic-encryption-based
computation. We first describe the implementation environment and
cryptographic parameters, then measure the running time and communication
of our protocol across a range of input sizes and network settings and
compare it against existing PJC protocols. Finally, we break down the
protocol's per-stage cost to identify the main contributors to overall
performance.



\subsection{Implementation and Experimental Setup}
\label{sec:experimental-setup}

\paragraph{Experiment environments.}
Our local measurements were run on Ubuntu 24.04.4 LTS with an AMD Ryzen
Threadripper 9970X 32-Core CPU and 64\,GB of RAM. We set the protocol
thread count to one for each party and used TCP over localhost. Linux
\textsf{tc netem} limited aggregate loopback egress to 10\,Gbps for LAN,
100\,Mbps for WAN100, and 10\,Mbps for WAN10. The target RTT was
0.2\,ms for LAN and 80\,ms for both WAN settings; half of the target RTT
was applied as delay on each loopback traversal. These are shared
aggregate-rate settings. Across the reported runs, the observed mean RTT
was 0.220--0.250\,ms for LAN, 80.064--80.094\,ms for WAN100, and
80.216--80.246\,ms for WAN10. No qdisc packet drops were observed.
Each table entry is from one execution of that configuration, with no
additional repetitions or separate boundary-test sweeps. Reported total
costs include setup; communication uses decimal MB, i.e., $10^6$ bytes.
The externally sourced values from \cite{CHI+23} retain their original
measurement conditions.

\medskip
\paragraph{Implementation details.}
We implemented our protocol in C++. AF-CPSI builds on \textsf{VOLE-PSI}~\cite{volepsi},
and we use \textsf{libOTe}~\cite{libOTe} for OT and B2A-OT, and we implement the homomorphic-encryption-based PJC
computation using the BGV~\cite{BGV14} scheme provided by Microsoft SEAL~\cite{SEAL}.

\medskip
\paragraph{Parameters.}
We use parameter settings targeting 128-bit computational security
and 40-bit statistical security, with Microsoft SEAL's \texttt{tc128}
parameter validation enabled. In the modular experiment, Ours uses
16-bit unsigned payloads, polynomial modulus degree $N=4096$, and
plaintext prime $p=4294475777\approx 2^{32}$. Its coefficient-modulus
bit sizes are $(48,36,18)$: the initial data-level ciphertext modulus
has approximately 84 bits, and the key modulus has 102 bits including
the 18-bit special prime. Homomorphic evaluation has multiplicative depth one: each input
ciphertext is multiplied by a plaintext, the results are accumulated,
and the accumulated ciphertext is modulus-switched.
Batching packs up to $N=4096$ values into a ciphertext.
In the exact integer inner-product setting (\Cref{sec:integer-ip}),
Ours uses unsigned 32-bit payloads, $N=8192$, and the two 43-bit plaintext
primes $p_1=8796092858369$ and $p_2=8796092792833$. Coefficient-modulus
bit sizes $(60,50,30)$ give an initial data-level modulus
$q\approx 2^{110}$ and a 140-bit key modulus including the 30-bit special
prime. The CRT product $P=p_1p_2$ has 86 bits and satisfies
$P>2^{20}(2^{32}-1)^2$, covering the full result range for the measured
set sizes. Table~\ref{tab:experimental-parameters} summarizes these
parameters.


\begin{table}[!htbp]
\centering
\caption{Cryptographic parameters used in our experiments.}
\label{tab:experimental-parameters}
\small
\begin{tabular}{@{}p{0.53\linewidth}p{0.41\linewidth}@{}}
\toprule
\textbf{Parameter} & \textbf{Setting} \\
\midrule
HE scheme & BGV \\
Modular payload width & 16 bits \\
Polynomial modulus degree $N$ & \(4096\) \\
Plaintext modulus $p$ & \(\approx 2^{32}\) \\
Ciphertext modulus $q$ & \(\approx 2^{84}\) \\
\midrule
\multicolumn{2}{l}{\textit{Integer inner product setting (\Cref{sec:integer-ip})}} \\
Polynomial modulus degree $N$ & $8192$ \\
Integer payload width & 32 bits \\
Plaintext moduli (RNS) & \(p_1,p_2\approx 2^{43}\) \\
Ciphertext modulus $q$ (data level) & \(\approx 2^{110}\) \\
Coefficient-modulus bits & $(60,50,30)$ \\
Key-modulus bit count & 140 (including special prime) \\
\bottomrule
\end{tabular}
\end{table}


\subsection{Comparison: Modular Inner Product}

\label{sec:performance-comparison}
This section reports the running time and communication of
\cite{CHI+23}, \cite{KLS26}, and Ours. The local experiments use the
existing modular implementations and input harnesses without changes.
Ours uses 16-bit unsigned payloads over $\mathbb{Z}_p$, with
$p=4294475777\approx 2^{32}$; the plaintext-modulus size does not denote
the payload width. The preserved \cite{KLS26} harness uses 32-bit integer
storage over the ring $\mathbb{Z}/(2^{32})$, assigning sender value $i$
and receiver value $-i$ and randomly retaining approximately half the
matches. Ours uses its existing deterministic payload generator and
exactly $n/2$ matches. Thus the table describes these respective
implementations and input settings; the payload distributions and
arithmetic domains differ. Since \cite{CHI+23} has no public
implementation, we retain its published results as reference values;
its reported payload width is $\ell=23$.



\begin{table}[!htbp]
\centering
\caption{Modular PJC performance under the revised network settings and the respective input settings described in the text.}
\label{tab:pjc-comparison}

\setlength{\tabcolsep}{2.2pt}
\renewcommand{\arraystretch}{1.05}
\footnotesize

\begin{tabular}{c l c c c c}
\toprule

\multirow{2}{*}{$n$}
&
\multirow{2}{*}{\textbf{Protocol}}
&
\multirow{2}{*}{\textbf{Comm. (MB)}}
&
\multicolumn{3}{c}{\textbf{Runtime (s)}}
\\

\cmidrule(lr){4-6}

&
&
&
\textbf{LAN}
&
\textbf{WAN}
&
\textbf{WAN}
\\[-2pt]

&
&
&
&
\textbf{100 Mbps}
&
\textbf{10 Mbps}
\\

\midrule

\multirow[t]{3}{*}{$2^{16}$}
& \cite{CHI+23}$^\dagger$& 17.2 & 25.92 & -- & -- \\
& \cite{KLS26} & 21.374 & 1.735 & 6.738 & 21.461 \\
\cmidrule(lr){2-6}%
& \textbf{Ours} & 14.150 & 0.750 & 4.567 & 14.118 \\

\midrule

\multirow[t]{3}{*}{$2^{20}$}
& \cite{CHI+23}$^\dagger$& 275.2 & 424.4 & -- & -- \\
& \cite{KLS26} & 330.100 & 28.112 & 60.469 & 302.270 \\
\cmidrule(lr){2-6}
& \textbf{Ours} & 218.010 & 12.426 & 32.710 & 188.588 \\

\bottomrule
\end{tabular}

\vspace{2pt}
\begin{minipage}{0.95\linewidth}
\footnotesize
$^\dagger$ \cite{CHI+23} has no public implementation, so we report the LAN numbers from the original paper.
\end{minipage}

\end{table}

Table~\ref{tab:pjc-comparison} reports lower observed total runtime
and communication for Ours than for the preserved \cite{KLS26}
implementation in all measured configurations. At $n=2^{16}$, Ours
communicates 14.150\,MB versus 21.374\,MB for \cite{KLS26}, giving a
communication ratio of $1.51$. The LAN times are 0.750\,s and
1.735\,s, respectively, a measured runtime ratio of $2.31$.

At $n=2^{20}$, Ours communicates 218.010\,MB versus
330.100\,MB for \cite{KLS26}, again a ratio of approximately $1.51$.
The LAN times are 12.426\,s and 28.112\,s, a runtime ratio of $2.26$.
Under WAN100 and WAN10, Ours takes 32.710\,s and 188.588\,s, compared
with 60.469\,s and 302.270\,s for \cite{KLS26}; the respective runtime
ratios are $1.85$ and $1.60$. These ratios are computed from the
unrounded, single-execution measurements under the stated input settings.

The \cite{CHI+23} rows retain the original paper's communication
and LAN runtime values. Its 23-bit payload setting and original
experimental environment differ from the local configurations, so
these entries serve as external reference points. The exact integer
comparison over unsigned 32-bit payloads is reported in
\Cref{sec:integer-ip-exp}.


\paragraph{Breakdown.}

We measured the per-stage running time and communication
over LAN for set size $2^{20}$, with three stages
-- AF-CPSI, B2A-OT, and homomorphic-encryption-based PJC computation --
and report the results in Table~\ref{tab:pjc-breakdown}.

This breakdown comes from an independent single execution at
$n=2^{20}$ under LAN 10\,Gbps with target RTT 0.2\,ms. Its total
runtime therefore differs from the separate execution used in
Table~\ref{tab:pjc-comparison}.

\begin{table}[!htbp]
\centering
\caption{Breakdown of our PJC protocol for $n=2^{20}$ over LAN 10\,Gbps with target RTT 0.2\,ms (independent single execution).}
\label{tab:pjc-breakdown}
\small
\setlength{\tabcolsep}{4pt}
\begin{tabular}{lcc}
\toprule
\textbf{Stage}
&
\textbf{Runtime (s)}
&
\textbf{Comm. (MB)}
\\
\midrule
AF-CPSI & 10.884 & 172.956 \\
B2A-OT & 0.347 & 22.769 \\
HE Computation & 0.261 & 22.285 \\
\midrule
\textbf{Total} & 11.518 & 218.010 \\
\bottomrule
\end{tabular}
\end{table}

In Table~\ref{tab:pjc-breakdown}, AF-CPSI accounts for approximately
94.50\% of the measured total runtime and 79.33\% of communication.
The total of 11.518\,s includes share addition and other overhead in
addition to the three displayed main stages. Share addition takes
0.019\,s; the total exceeds the sum of the three displayed stages by
approximately 0.025\,s. The separately logged VOLE calibration is a
diagnostic run and is excluded from protocol totals. 

\subsection{Comparison: Integer Inner Product}
\label{sec:integer-ip-exp}
We measured the setting of \Cref{sec:integer-ip}, i.e., the exact integer
inner product over 32-bit payloads, under LAN and two WAN settings.
 
For this experiment, we adapt the original 32-bit-ring implementation
of \cite{KLS26} to actual arithmetic in $\mathbb{Z}/(2^{84})$.
The OKVS value shares and OT arithmetic messages use 11-byte storage
with reduction to 84 bits. MultShare uses 32 OT choice bits per payload,
matching the unsigned 32-bit input width. The implementation performs
arithmetic addition, subtraction, multiplication, and shifts, computes
the inner product, and checks it against a direct unsigned 128-bit
calculation in every measured execution. These measurements describe
our adaptation of \cite{KLS26}. Both implementations use full-width
unsigned 32-bit payloads with the same payload formulas and $n/2$
matches. The result is bounded by $n(2^{32}-1)^2<2^{84}$ for
$n\leq 2^{20}$. Ours reconstructs the same exact integer using the two
43-bit CRT primes described above; its CRT arithmetic space has 86 bits.

\begin{table}[!htbp]
\centering
\caption{Exact integer inner product over unsigned 32-bit payloads under the revised network settings.}
\label{tab:pjc-long}
\setlength{\tabcolsep}{2.2pt}
\renewcommand{\arraystretch}{1.05}
\footnotesize
\begin{tabular}{c l c c c c}
\toprule
\multirow{2}{*}{$n$} & \multirow{2}{*}{\textbf{Protocol}} & \multirow{2}{*}{\textbf{Comm. (MB)}} & \multicolumn{3}{c}{\textbf{Runtime (s)}} \\
\cmidrule(lr){4-6}
& & & \textbf{LAN} & \textbf{100 Mbps} & \textbf{10 Mbps} \\
\midrule
\multirow[t]{2}{*}{$2^{16}$}
& \cite{KLS26} & 44.711 & 2.498 & 9.760 & 41.078 \\
& \textbf{Ours} & 22.090 & 1.156 & 5.671 & 20.665 \\
\midrule
\multirow[t]{2}{*}{$2^{20}$}
& \cite{KLS26} & 708.059 & 40.039 & 103.678 & 618.093 \\
& \textbf{Ours} & 340.694 & 15.869 & 43.972 & 291.313 \\
\bottomrule
\end{tabular}
\vspace{2pt}
\begin{minipage}{0.95\linewidth}
\footnotesize
\cite{KLS26} denotes our adaptation using actual 84-bit arithmetic
and 32-bit OT choices; every measured execution computes and verifies
the integer result. Ours uses two 43-bit CRT primes. Both spaces cover
the exact result range for $n\leq 2^{20}$.
\cite{CHI+23} has no public implementation and reports only $\ell=23$.
\end{minipage}
\end{table}

The modular and integer experiments use different payload and HE
settings. For Ours, the payload width changes from 16 to 32 bits,
the residue count changes from one to two, each plaintext prime changes
from approximately 32 to 43 bits, and $N$ changes from 4096 to 8192.
At $n=2^{20}$, Ours communicates 218.010\,MB in the modular setting
and 340.694\,MB in the integer setting, an increase by a factor of
$1.56$. The corresponding \cite{KLS26} totals are 330.100\,MB and
708.059\,MB, a factor of $2.14$. These measured configuration-to-configuration
ratios combine the changes in input and arithmetic settings and do not
isolate an asymptotic dependence on arithmetic-space width.

For the exact integer task, the communication ratios
\cite{KLS26}/Ours are $2.02$ at $n=2^{16}$ and $2.08$ at
$n=2^{20}$. The respective LAN runtime ratios are $2.16$ and $2.52$.
At $n=2^{20}$, the WAN100 times are 103.678\,s for \cite{KLS26}
and 43.972\,s for Ours, a ratio of $2.36$; the WAN10 times are
618.093\,s and 291.313\,s, a ratio of $2.12$. These values describe
the observed single-execution performance in each network setting.
